\chapter{Генерация с извлечением информации: лексический, гибридный и двухэтапный поиск}\label{ch:rag}

Генерация с извлечением информации (Retrieval-Augmented Generation, RAG) объединяет извлечение релевантных документов из внешнего хранилища с генеративными возможностями больших языковых моделей.
В обзоре литературы (Раздел~\ref{sec:rag}) систематизированы основные парадигмы RAG; настоящая глава фокусируется на построении конкретного поискового стека --- лексического, плотного, гибридного и двухэтапного поиска --- и раскрывает назначение RAG как механизма, обеспечивающего достоверность и актуальность ответов.
Описываемый стек образует основу безопасного индекса, используемого в конвейере обеспечения безопасности (Глава~\ref{ch:safety}).

%% =================================================================
\section{Назначение и преимущества RAG}\label{sec:rag_motivation}
%% =================================================================

Знания большой языковой модели хранятся в её параметрах (параметрическая память) и фиксируются на момент завершения обучения.
Это порождает два фундаментальных ограничения: модель склонна к генерации правдоподобных, но фактически неверных утверждений (галлюцинаций), а её знания устаревают и не могут быть обновлены без переобучения.
Основополагающая работа Lewis et al.~\cite{2005.11401} предложила дополнить параметрическую память непараметрической --- внешним индексом документов, из которого релевантная информация извлекается на этапе инференса и передаётся модели в качестве контекста.
Данный подход обеспечивает ряд принципиальных преимуществ.

\textbf{Снижение галлюцинаций.} Генерация, заземлённая (grounded) на извлечённых документах, опирается на проверяемые источники, а не исключительно на параметрическую память модели.
Это существенно снижает частоту фактических ошибок и позволяет сопровождать ответ ссылками на источники (атрибуцией), обеспечивая верифицируемость.

\textbf{Отсутствие необходимости переобучения при появлении новых данных.} Обновление знаний системы сводится к обновлению внешнего индекса, а не весов модели.
Добавление, изменение или удаление документов происходит без дорогостоящего дообучения, что обеспечивает актуальность системы при динамично меняющихся данных и возможность оперативного отзыва недостоверной информации.

\textbf{Доменная адаптация и контролируемость.} Подключение специализированного индекса позволяет адаптировать модель к узкой предметной области без изменения её параметров, а ограничение извлечения определённым набором источников --- контролировать допустимое информационное пространство, что особенно важно для систем обеспечения безопасности.

%% =================================================================
\section{Лексический поиск: BM25}\label{sec:rag_bm25}
%% =================================================================

Лексический (разреженный) поиск оценивает релевантность по совпадению термов запроса и документа.
Классической функцией ранжирования является BM25~\cite{robertson2009bm25}, основанная на вероятностной модели релевантности.
Оценка релевантности документа $d$ запросу $q$ вычисляется как сумма вкладов входящих в запрос термов $t$:
\begin{equation}
\label{eq:bm25}
    \mathrm{BM25}(q, d) = \sum_{t \in q} \mathrm{IDF}(t) \cdot \frac{f(t, d) \cdot (k_1 + 1)}{f(t, d) + k_1 \cdot \left(1 - b + b \cdot \dfrac{|d|}{\mathrm{avgdl}}\right)},
\end{equation}
где $f(t, d)$ --- частота терма $t$ в документе $d$, $|d|$ --- длина документа, $\mathrm{avgdl}$ --- средняя длина документа в коллекции, а $k_1$ и $b$ --- гиперпараметры, регулирующие насыщение по частоте терма и степень нормализации по длине документа (типичные значения $k_1 \in [1.2, 2.0]$, $b = 0.75$).
Обратная документная частота $\mathrm{IDF}(t)$ снижает вклад часто встречающихся термов:
\begin{equation}
\label{eq:bm25_idf}
    \mathrm{IDF}(t) = \ln\frac{N - n(t) + 0.5}{n(t) + 0.5},
\end{equation}
где $N$ --- общее число документов в коллекции, $n(t)$ --- число документов, содержащих терм $t$.
Лексический поиск точно обрабатывает редкие термины, имена собственные и точные совпадения, эффективно реализуется через инвертированный индекс, однако не учитывает синонимию и семантическую близость при отсутствии лексического пересечения.

%% =================================================================
\section{Плотный поиск}\label{sec:rag_dense}
%% =================================================================

Плотный поиск преодолевает ограничения лексического подхода, оценивая семантическую близость в непрерывном векторном пространстве (Уравнение~\ref{eq:rr_biencoder}).
Dense Passage Retrieval~\cite{2004.04906} продемонстрировал, что плотные представления, полученные двухэнкодерной архитектурой, существенно превосходят разреженные методы на задачах поиска пассажей.
В настоящей работе в качестве модели плотного поиска используется GigaEmbeddings (Глава~\ref{ch:gigaembeddings}), а поиск реализуется через приближённый поиск ближайших соседей в векторном индексе.
Плотный поиск улавливает семантическую близость даже при отсутствии лексического пересечения, однако может уступать лексическому на запросах с редкими терминами и точными формулировками.

%% =================================================================
\section{Гибридный поиск}\label{sec:rag_hybrid}
%% =================================================================

Лексический и плотный подходы дополняют друг друга: первый силён в точных совпадениях, второй --- в семантическом обобщении.
Гибридный поиск (Рисунок~\ref{fig:hybrid_search}) объединяет их результаты, повышая полноту и устойчивость выдачи.

\begin{figure}[htbp]
    \centering
    \resizebox{0.92\textwidth}{!}{
        \begin{tikzpicture}[
    font=\sffamily,
    node distance=1.2cm and 1.9cm,
    box/.style={
        draw,
        rounded corners,
        minimum width=3.0cm,
        minimum height=1.1cm,
        align=center,
        thick
    },
    db/.style={
        draw,
        cylinder,
        shape border rotate=90,
        aspect=0.3,
        minimum width=2.0cm,
        minimum height=1.4cm,
        align=center,
        thick
    },
    fuse/.style={
        draw,
        diamond,
        aspect=2,
        minimum width=3.0cm,
        align=center,
        thick
    },
    arr/.style={-{Stealth[length=3mm]}, thick}
]
    %% Запрос
    \node[box] (query) {Запрос};
    %% Две ветви поиска
    \node[box, above right=0.6cm and 2.0cm of query] (sparse) {Лексический поиск\\BM25};
    \node[box, below right=0.6cm and 2.0cm of query] (dense) {Плотный поиск\\(GigaEmbeddings)};
    \node[db, above=of sparse] (invindex) {Инвертир.\\индекс};
    \node[db, below=of dense] (vecindex) {Векторный\\индекс};
    %% Слияние
    \node[fuse] (rrf) at ($(sparse)!0.5!(dense) + (3.4cm,0)$) {Слияние\\рангов (RRF)};
    %% Результат
    \node[box, right=of rrf] (topk) {Top-$k$\\документов};

    %% Стрелки
    \draw[arr] (query) -- (sparse);
    \draw[arr] (query) -- (dense);
    \draw[arr] (invindex) -- (sparse);
    \draw[arr] (vecindex) -- (dense);
    \draw[arr] (sparse) -- (rrf);
    \draw[arr] (dense) -- (rrf);
    \draw[arr] (rrf) -- (topk);
\end{tikzpicture}

    }
    \caption{Схема гибридного поиска: результаты лексического (BM25) и плотного (GigaEmbeddings) поиска объединяются слиянием рангов.}
    \label{fig:hybrid_search}
\end{figure}

Распространённым методом объединения, не требующим калибровки разнородных оценок, является слияние по обратным рангам (Reciprocal Rank Fusion, RRF)~\cite{cormack2009rrf}.
Каждому документу $d$ присваивается агрегированная оценка на основе его позиций (рангов) в выдачах отдельных систем:
\begin{equation}
\label{eq:rrf}
    \mathrm{RRF}(d) = \sum_{r \in \mathcal{R}} \frac{1}{\kappa + \mathrm{rank}_r(d)},
\end{equation}
где $\mathcal{R}$ --- множество объединяемых систем поиска (лексической и плотной), $\mathrm{rank}_r(d)$ --- ранг документа $d$ в выдаче системы $r$, а $\kappa$ --- сглаживающая константа (типичное значение $\kappa = 60$), снижающая влияние документов с высокими рангами.
Альтернативный подход --- выпуклая комбинация нормированных оценок:
\begin{equation}
\label{eq:hybrid_convex}
    s_{\text{hybrid}}(q, d) = \alpha \cdot \tilde{s}_{\text{dense}}(q, d) + (1 - \alpha) \cdot \tilde{s}_{\text{sparse}}(q, d),
\end{equation}
где $\tilde{s}$ обозначают нормированные оценки соответствующих систем, а $\alpha \in [0, 1]$ регулирует баланс между плотным и лексическим поиском и подбирается на валидационной выборке.
Слияние по обратным рангам, как правило, предпочтительнее в силу устойчивости к различию масштабов оценок разнородных систем.

%% =================================================================
\section{Двухэтапный поиск}\label{sec:rag_two_stage}
%% =================================================================

Полный поисковый стек организован в виде двухэтапного конвейера (Раздел~\ref{sec:rr_two_stage}).
На этапе отбора кандидатов применяется лексический, плотный или гибридный поиск, формирующий множество из top-$k$ документов с высокой полнотой.
На этапе уточнения это множество переупорядочивается кросс-энкодером (Глава~\ref{ch:reranking}), обеспечивающим высокую точность за счёт совместного кодирования пар <<запрос--документ>>.
Полученное на втором этапе множество top-$n$ наиболее релевантных документов передаётся языковой модели в качестве контекста для генерации заземлённого ответа.
Такая организация объединяет преимущества всех компонентов: масштабируемость отбора, полноту гибридного поиска и точность переранжирования.

%% =================================================================
\section{Связь с безопасным индексом}\label{sec:rag_safe_index}
%% =================================================================

Описанный двухэтапный поисковый стек служит основой безопасного индекса (safe index), используемого в конвейере обеспечения безопасности (Глава~\ref{ch:safety}).
Безопасный индекс представляет собой курируемое хранилище проверенных и актуальных документов, доступ к которому осуществляется через инструмент безопасного поиска (safe-search-tool).
В отличие от поиска по произвольному веб-индексу, поиск ограничивается заранее отфильтрованным информационным пространством, что обеспечивает безопасной модели актуальный и достоверный взгляд на предметную область без риска извлечения вредоносного или недостоверного контента.
Тем самым безопасный индекс представляет собой не просто отфильтрованный корпус, а механизм ограничения информационного пространства модели: контролируя множество доступных источников, он становится частью защитного контура наряду с фильтрацией входов и выходов (см. ключевую гипотезу работы во Введении).
Тем самым поисковый стек, развиваемый в части диссертации, посвящённой эмбеддингам (Глава~\ref{ch:gigaembeddings}), переранжированию (Глава~\ref{ch:reranking}) и RAG (настоящая глава), непосредственно интегрируется в систему обеспечения безопасности больших языковых моделей.

%% =================================================================
\section{Выводы}\label{sec:rag_conclusions}
%% =================================================================

В настоящей главе рассмотрено назначение генерации с извлечением информации как механизма, снижающего галлюцинации и устраняющего необходимость переобучения при появлении новых данных за счёт обновления внешнего индекса.
Представлен поисковый стек, включающий лексический поиск BM25, плотный поиск на основе GigaEmbeddings, гибридный поиск со слиянием по обратным рангам и двухэтапную схему с переранжированием.
Показано, что данный стек образует основу безопасного индекса, связывая поисковую и защитную части настоящего исследования в единую систему, рассматриваемую в Главе~\ref{ch:safety}.
Сквозная экспериментальная оценка собранного поискового стека (recall, nDCG@$k$, MRR@$k$, задержка) на русскоязычных коллекциях составляет направление дальнейшей работы.
